% !TEX program = xelatex

\documentclass{resume}
\usepackage{zh_CN-Adobefonts_external} % Simplified Chinese Support using external fonts (./fonts/zh_CN-Adobe/)
% \usepackage{zh_CN-Adobefonts_internal} % Simplified Chinese Support using system fonts
\usepackage{hyperref}
\begin{document}
\pagenumbering{gobble} % suppress displaying page number

\name{Zheng Sun}

\basicInfo{
  zhsun@aip.de |
  +49 1623062086 | 
  \href{https://scholar.google.com/citations?user=lBU4pfsAAAAJ&hl=zh-CN}{Google Scholar} |
  \href{https://orcid.org/0000-0001-5657-7587}{Orcid}
}


\section{Research Interests}
\begin{itemize}
\item \textbf{Solar and Stellar Eruptions}
\item \textbf{Star-planet Interaction}
\end{itemize}

\section{Education}
\datedsubsection{\textbf{Peking University (PKU)}, Beijing, China}{Sep. 2022 -- Jan. 2028}
\textbf{Ph.D. Candidate, Space Physics}, School of Earth and Space Sciences
\\Supervisor: \href{https://scholar.google.com/citations?user=uoTD4VUAAAAJ&hl=zh-CN&oi=ao}{ Prof. Hui Tian }

\datedsubsection{\textbf{Leibniz Institute for Astrophysics Potsdam (AIP)}, Potsdam, Germany}{Nov. 2024 -- Nov. 2026}
\textbf{Visiting Student, Stellar Physics and Exoplanets}
\\Supervisor: \href{https://pages.aip.de/jalvarado/}{ Dr. Julián D. Alvarado-Gómez }




\datedsubsection{\textbf{China University of Geosciences (CUGB)}, Beijing, China}{Sep. 2018 -- Jun. 2022}
\textbf{Bachelor, Geophysics}, School of Geophysics and Information Technology


\section{Research Experience}

\datedsubsection{\textbf{Numerical Simulations of Exoplanetary Space Weather}}{Nov. 2024 -- Present}
\role{Ph.D. Project}{Supervisor: Prof. Hui Tian \& Dr. Julián D. Alvarado-Gómez}
\begin{itemize}
\item Using the Space Weather Modeling Framework (SWMF) to simulate stellar wind and coronal mass ejections (CMEs), exploring the solar--stellar connection and the impact of CMEs on exoplanets
\item Simulating energetic particle environments in exoplanetary systems
\end{itemize}

\datedsubsection{\textbf{Solar Active Region Evolution Associated with Eruptions}}{Sep. 2022 -- Present}
\role{Ph.D. Project}{Supervisor: Prof. Hui Tian \& Dr. Ting Li}
\begin{itemize}
\item Investigated differences between eruptive and confined solar flares during their pre-eruptive phases using SDO/AIA and other EUV imaging observations.
\item Studied the evolution of magnetic helicity and current helicity in solar active regions using SDO/HMI magnetograms and examined their relationship with solar eruptions.
\end{itemize}
\section{Observing Experience}

\datedsubsection{\textbf{Solar Flare Observations with VTT}}{Jul. 3--15, 2026}
\role{Observer}{Teide Observatory, Tenerife, Spain}

\datedsubsection{\textbf{Stellar Spectropolarimetric Observations with HARPSpol (PHAME Survey)}}{Dec. 20--28, 2025}
\role{Observer}{La Silla Observatory, Chile}



% Reference Test
%\datedsubsection{\textbf{Paper Title\cite{zaharia2012resilient}}}{May. 2015}
%An xxx optimized for xxx\cite{verma2015large}
%\begin{itemize}
%  \item main contribution
%\end{itemize}
\newpage
\section{Publications}
\textbf{First-author publications:}
\begin{enumerate}
\item \textbf{Sun, Z.}, Alvarado-Gómez, J.D., Tian, H., et al. (2025),
``A Stellar Coronal Mass Ejection Candidate Detected on the Young Solar-Type Star DS Tucana A with VLT/ESPRESSO'',
submitted to \textit{The Astrophysical Journal}.
\item \textbf{Sun, Z.}, Alvarado-Gómez, J.D., Tian, H., et al. (2025),
``How Magnetic Field Strength Governs Stellar Coronal Mass Ejection Dynamics: A Parametric Study'',
submitted to \textit{The Astrophysical Journal}.
\item \textbf{Sun, Z.}, Alvarado-Gómez, J.D., Tian, H., et al. (2025),
``Intense but Harmless: Space Weather around an M dwarf'',
submitted to \textit{The Astrophysical Journal}.
\item \textbf{Sun, Z.}, Alvarado-Gómez, J.D., Tian, H., et al. (2025),
``Fine-scale downflows above flare ribbons captured by Solar Orbiter/EUI'',
accepted by \href{https://arxiv.org/abs/2606.19335}{\textit{Astronomy \& Astrophysics}}.
\item \textbf{Z. Sun}, T. Li, X. Bian, et al. (2025). ``Current Helicity in Response to Coronal Mass Ejections". \href{https://iopscience.iop.org/article/10.3847/1538-4357/adf18b/meta}{\textit{The Astrophysical Jounal, 990(1), 45.}}
\item \textbf{Z. Sun}, H. Tian, T. Li, et al. (2024). ``Fast Downflows Observed during a Polar Crown Filament Eruption". \href{https://iopscience.iop.org/article/10.3847/1538-4357/ad738d}{\textit{The Astrophysical Journal, 974(2), 205.}}
\item \textbf{Z. Sun}, T. Li, Y. Hou, et al. (2024). ``The Origin of an Intense Geomagnetic Storm on 2023 December 1st: Successive Buildup and Eruption of Multiple Magnetic Flux Ropes". \href{https://link.springer.com/article/10.1007/s11207-024-02329-4}{\textit{Solar Physics, 299(6), 1-22.}}
\item \textbf{Z. Sun}, T. Li, Q. Wang, et al. (2024). ``Magnetic Helicity Evolution during Active Region Emergence and Subsequent Flare Productivity". \href{https://www.aanda.org/articles/aa/abs/2024/06/aa48734-23/aa48734-23.html}{\textit{Astronomy \& Astrophysics, 686, A148.}}
\item \textbf{Z. Sun}, T. Li, H. Tian, et al. (2023). ``Observation of Two Splitting Processes in a Partial Filament Eruption on the Sun: the Role of Breakout Reconnection". \href{https://iopscience.iop.org/article/10.3847/1538-4357/ace5b1/meta}{\textit{The Astrophysical Journal, 953(2), 148.}}
\item \textbf{Z. Sun}, H. Tian, P. F. Chen, et al. (2022). ``Cross-loop Propagation of a Quasiperiodic Extreme-Ultraviolet Wave Train Triggered by Successive Stretching of Magnetic Field Structures during a Solar Eruption". \href{https://iopscience.iop.org/article/10.3847/2041-8213/ac9aff/meta}{\textit{The Astrophysical Journal Letters, 939(2), L18.}}



\end{enumerate}


\textbf{Co-authored publications:}
\begin{enumerate}[start=11]
\item Z. Wu, T. Van Doorsselaere, J. He, …, \& \textbf{Z. Sun}. (2026). ``Imaging Magnetically Driven Astrosphere". Accepted by \textit{The Astrophysical Journal Letters.}
\item X. Duan, T. Li, Y. Cui, … \& \textbf{Z. Sun}. (2026). ``Solar Energetic Particle Events and Associated Type II Radio Bursts from Different Source Regions". \href{https://iopscience.iop.org/article/10.3847/1538-4357/ae63bb}{\textit{The Astrophysical Journal, 1003(1), 53.}}
\item F. Qiao, H. Li, J. Wang, Y. Duan, \textbf{Z. Sun}, et al. (2025). ``Fine Structure and Formation Mechanism of a Sunspot Bipolar Light Bridge in NOAA AR 13663". \href{https://iopscience.iop.org/article/10.3847/1538-4357/ae4bd9}{\textit{The Astrophysical Journal, 1000(2), 205.}}
\item Y. Duan, X. Yan, J. Hong, …, \textbf{Z. Sun}, et al. (2025). ``Resolving Interchange Reconnection Dynamics in a Fan-Spine-like Topology Observed by Solar Orbiter". \href{https://www.aanda.org/articles/aa/abs/2026/02/aa57158-25/aa57158-25.html}{\textit{Astronomy \& Astrophysics, 706, A1.}}
\item X. Duan, T. Li, Y. Hou, …, \textbf{Z. Sun}, et al. (2025). ``A Statistical Analysis of Magnetic Parameters in Solar Source Regions of Halo Coronal Mass Ejections with and without Solar Energetic Particle Events". \href{https://iopscience.iop.org/article/10.3847/1538-4357/ade9af}{\textit{The Astrophysical Journal, 988(2), 265.}}
\item Y. Chen, H. Lu, H. Tian, …, \textbf{Z. Sun}, et al. (2025). ``Asymmetries in Fe IX 17.1 nm and Fe XII 19.5 nm Line Profiles as Possible Signatures of Obscuration Dimming in Sun-as-a-star Spectra". \href{https://iopscience.iop.org/article/10.3847/2041-8213/ade685}{\textit{The Astrophysical Journal Letters, 987(1), L22.}}
\item Y. Zhang, T. Li, W. Teng, …, \textbf{Z. Sun}, et al. (2025). ``Quasiperiodic Slow-propagating Extreme-ultraviolet “Wave” Trains after the Filament Eruption". \href{https://iopscience.iop.org/article/10.3847/2041-8213/add33b}{\textit{The Astrophysical Journal Letters, 985(1), L1.}}
\item L. Chan, X. Bai, Y. Feng, … \& \textbf{Z. Sun} (2025). ``The Detectability of Coronal Mass Ejections in the Low Corona Using Multi-slit Extreme-ultraviolet Spectroscopy". \href{https://iopscience.iop.org/article/10.3847/1538-4357/adc10a}{\textit{The Astrophysical Journal, 984(2), 141.}}
\item Y. Zhang, T. Li, Y. Hou, X. Duan, \textbf{Z. Sun}, et al. (2025)  ``Quasiperiodic Super-Alfvénic Slippage along Flare Ribbons Observed by the Interface Region Imaging Spectrograph". \href{https://iopscience.iop.org/article/10.3847/2041-8213/adbb6e}{\textit{The Astrophysical Journal Letters, 982(1), L9.}}
\item Y. Duan, H. Chen, Z. Hou, \textbf{Z. Sun}, et al. (2025). ``Ubiquitous Small-scale Extreme-ultraviolet Upflow-like Events above Network Regions Observed by the Solar Orbiter/Extreme Ultraviolet Imager". \href{https://iopscience.iop.org/article/10.3847/1538-4357/ada556}{\textit{The Astrophysical Journal, 979(2), 195.}}
\item H. Lu, H. Tian, L. Zhang, … \& \textbf{Z. Sun} (2025). ``An Extreme Stellar Prominence Eruption Observed by LAMOST Time-Domain Spectroscopy". \href{https://iopscience.iop.org/article/10.3847/2041-8213/ad93cc/meta}{\textit{The Astrophysical Journal Letters, 978(2), L32.}}
\item F. Qiao, L. Li, H. Tian, … \& \textbf{Z. Sun} (2024). ``Three Types of Solar Coronal Rain during Magnetic Reconnection between Open and Closed Magnetic Structures". \href{https://iopscience.iop.org/article/10.3847/1538-4357/ad6770}{\textit{The Astrophysical Journal, 973(1), 57.}}
\item Z. Yang, H. Tian, Y. Zhu, … \& \textbf{Z. Sun} (2024). ``Is It Possible to Detect Coronal Mass Ejections on Solar-Type Stars through Extreme-Ultraviolet Spectral Observations?". \href{https://iopscience.iop.org/article/10.3847/1538-4357/ad2a44}{\textit{The Astrophysical Journal, 966(1), 24.}}
\item Y. Duan, H. Tian, H. Chen, Y. Shen, \textbf{Z. Sun}, et al.  (2024). ``Formation of Fan-spine Magnetic Topologies through Flux Emergence and Subsequent Jet Production". \href{https://iopscience.iop.org/article/10.3847/2041-8213/ad24f3/meta}{\textit{The Astrophysical Journal Letters, 962(2), L38.}}
\item X. Bai, H. Tian, Y. Deng, …, \textbf{Z. Sun}, et al. (2023). ``The Solar Upper Transition Region Imager (SUTRI) onboard the SATech-01 Satellite". \href{https://iopscience.iop.org/article/10.1088/1674-4527/accc74/meta}{\textit{Research in Astronomy and Astrophysics}, 23(6), 065014.}

\end{enumerate}



\newpage
\section{Conferences}
\textbf{Contributed Talks:}
\begin{itemize}
\item IAU Symposium 408: Unraveling the joint lives of Stars and Exoplanets, Aug.17-21, Liege, Belgium, ``Space Weather around M dwarfs: Role of Coronal Mass Ejections"
\item 46th COSPAR Scientific Assembly, Florence, Italy, Aug.1-9, 2026, ``Space Weather around M dwarfs: Role of Coronal Mass Ejections"
\item European Astronomical Society Annual Meeting 2026, Jun.29-July.3, Lausanne, Switzerland, ``Space Weather around M dwarfs: Role of Coronal Mass Ejections"
\item EGU General Assembly 2026, Vienna, Austria, May.3-8, 2026, ``Current Helicity Reversal during Coronal Mass Ejections"
\item Berlin Early-career Space Research Conference 2025, Berlin, Germany, Oct.6-7, 2025, ``Space Weather around Moderately-rotating, Fully Convective M-dwarfs"
\item IAU Symposium 400: Solar and Stellar Multi-scale Activitity, Medellín, Colombia, Jul.21-25, 2025, ``Space Weather around Moderately-rotating, Fully Convective M-dwarfs"
\item European Astronomical Society Annual Meeting 2025, Cork, Ireland, Jun.23-27, 2025, ``Space Weather around Moderately-rotating, Fully Convective M-dwarfs"
\item 6th SWMF User Meeting, online, Mar.6-7, 2025, ``Space Weather around Moderately-rotating, Fully Convective M-dwarfs"
\item 4th Solar Activity Workshop, Weihai, China, Aug.12-15, 2024, ``Fast Downflows in a Polar Crown Filament Eruption"
\item 45th COSPAR Scientific Assembly, Busan, South Korea, Jul.13-21, 2024, ``Magnetic Helicity Evolution during Active Region Emergence and Subsequent Flare Productivity" 
\item ISSI-BJ International Team Meeting on Magnetohydrodynamic Wavetrains as a Tool for Probing the Solar Corona, Beijing, China, Dec.4-8, 2023, ``Cross-loop EUV Wave Train Triggered by Successive Stretching of Magnetic Field Structures" 
\item 1st Joint Conference of ASO-S and CHASE Missions, Wuxi, China, Oct.10-13, 2023, ``Two Splitting Processes in a Partial Filament Eruption"  
\item 4th National Digital Space Strategy Symposium and 13th National Space Weather Symposium, Kunming, China, Jul.18-22, 2023, ``Magnetic Helicity Evolution during Active Region Emergence and Subsequent Flare Productivity" 
\item National Solar Physics Meeting 2023, Chengdu, China, Jun.7-11, 2023, ``Cross-loop EUV Wave Train Triggered by Successive Stretching of Magnetic Field Structures"  
\item Workshop on Waves in the Solar and Stellar Atmosphere, online, Apr.4, 2023, ``Cross-loop EUV Wave Train Triggered by Successive Stretching of Magnetic Field Structures"  
\end{itemize}

\textbf{Posters:}
\begin{itemize}
\item Cool Stars 23, Tokyo, Japan, Jun.15-19, 2026, ``Space Weather around M dwarfs: Role of Coronal Mass Ejections"
\item Solar Orbiter Workshop 2026, Berlin, Germany, Mar.16-19, 2026, ``Fine-scale Downflows Above Flare Ribbons Captured by Solar Orbiter/EUI"
\item 17th European Solar Physics Meeting, Turin, Italy, Sep.9-13, 2024, ``Magnetic Helicity Evolution during Active Region Emergence and Subsequent Flare Productivity" 
\item 1st Sun and Space Weather Joint Conference: Extreme Solar Eruption Events and Their Interplanetary Processes, Altay, China, Aug.18-21, 2024, ``Fast Downflows in a Polar Crown Filament Eruption"
\item IAU Symposium 388: Solar and Stellar Coronal Mass Ejections, Krakow, Poland, May.5-10, 2024, ``Fast Downflows in a Polar Crown Filament Eruption" 
\end{itemize}



\textbf{Attendances:}
\begin{itemize}
\item Berlin-Potsdam Exoplanet Community: Host Stars, Atmospheres, Interiors and Geodynamics (EXCHANGE) Roundtable, Berlin, Germany, Nov.29, 2024, and Dec.5, 2025 and Potsdam, Germany, Mar.14, 2025
\item High-Energy Cosmic Rays and Solar/Interplanetary Physics Symposium, Beijing, China, Jan.19-21, 2024
\item Mini-workshop on CME Initiation, Nanjing, China, Nov.18, 2023
\item Major Eruptions and Space Weather in Solar Cycle 25 Project Meeting, Beijing, China, May.18, 2023
\item Workshop on Solar Atmospheric Magnetic Field Diagnosis and Related Physical Process Based on Spectral Imaging, Beijing, China, Feb.14-16, 2023
\end{itemize}



\newpage

\section{Professional Experience}

\datedsubsection{\textbf{Remote Operations of Solar Upper Transition Region Imager (SUTRI)}}{Nov. 2022 -- Sep. 2024}
\role{Operator}{National Astronomical Observatories, Chinese Academy of Sciences}

\datedsubsection{\textbf{Solar EUV Data Processing Pipeline Development for FengYun-3E}}{Sep. 2022 -- Jul. 2024}
\role{Contributor}{China Meteorological Administration}

\datedsubsection{\textbf{Geophysical Field Survey}}{Jun. 2021 -- Jul. 2021}
\role{Student Researcher}{Liujiang Basin, Qinhuangdao, China}


\section{Professional Service}

\begin{itemize}
\item Referee for \textit{Astronomische Nachrichten}, Mar. 2026
\item Referee for \textit{Monthly Notices of the Royal Astronomical Society} (MNRAS), Jun. 2025
\end{itemize}

\section{Honors and Awards}
\begin{itemize}
\item \datedline{\textit{National Scholarship, PKU}}{Dec. 2025}
\item \datedline{\textit{National Scholarship, PKU}}{Dec. 2024}
\item \datedline{\textit{Chinese Government Scholarship, China Scholarship Council}}{Jul. 2024}
\item \datedline{\textit{Ubiquant Scholarship, PKU}}{Dec. 2023}
\item \datedline{\textit{Outstanding Undergraduate Thesis, CUGB}}{Jun. 2022}
\item \datedline{\textit{Greenview Scholarship, CUGB}}{Dec. 2021}
\item \datedline{\textit{Silvercorp Scholarship, CUGB}}{Dec. 2020}
\item \datedline{\textit{China National Petroleum Corporation Scholarship, CUGB}}{Dec. 2019}

\end{itemize}


\section{Activities}
\begin{itemize}
\item \datedline{\textit{Third Place, Floorball Freshman Cup, PKU}}{Nov. 2024}
\item \datedline{\textit{Second Place, Volleyball PKU Cup, PKU}}{Jun. 2024}
\item \datedline{\textit{Top 8, Openfire MC Freestyle Battle, Beijing}}{Apr. 2023}
\item \datedline{\textit{Champion, Basketball Association Cup, CUGB}}{Oct. 2018}
\end{itemize}
%% Reference
%\newpage
%\bibliographystyle{IEEETran}
%\bibliography{mycite}


\section{Language Proficiency}
Chinese: Native Speaker
\\English: Fluent
\\German: Basic

\end{document}
